\subsection{GRAPHS AND CORRESPONDENCES}

DEFINITION 1. G is said to be a graph if every element of G is an ordered pair, i.e., if the relation

\[
(\forall z) (z \in G \Longrightarrow (z \text {   is   an   ordered   pair }))
\]

is true.

If G is a graph, the relation \((x, y) \in \mathbf{G}\) is expressed by saying that ``y corresponds to x under G''.

Let \(R \mid x, y \mid\) be a relation, where \(x\) and \(y\) are distinct letters. Let \(G\) be a letter, distinct from \(x\) and \(y\) , which does not appear in \(R\) . If the relation

\[
(\exists G) (G \text {   is   a   graph   and   } (\forall x) (\forall y) (R \Longleftrightarrow ((x, y) \in G)))
\]

is true, the relation R is said to have a graph (with respect to the letters x and y). The graph G is then unique by virtue of the axiom of extent, and is called the graph of R with respect to x and y.

Let \(\pmb{Z}\) be a letter, distinct from \(\pmb{x}\) and \(\pmb{y}\) , which does not appear in \(\pmb{R}\) . If the relation

\[
(\exists Z) (\forall x) (\forall y) (R \Longrightarrow ((x, y) \in Z))
\]

is true, then \(R\) has a graph; we may take this graph to be the set of ordered pairs \(z\) such that \(z \in Z\) and \(R \{ \text{pr}_1 z, \text{pr}_2 z \}\) ( \(z\) being a letter distinct from \(x, y, Z\) which does not appear in \(R\) ). This condition is satisfied if we know a term \(T\) , which does not contain either \(x\) or \(y\) , such that \(R \Longrightarrow ((x, y) \in T)\) is true.

PROPOSITION 1. Let G be a graph. There exists exactly one set A and exactly one set B with the following properties:

\begin{enumerate}
\def\labelenumi{(\arabic{enumi})}
\item
  the relation \((\exists y)((x, y) \in G)\) is equivalent to \(x \in A\) ;
\item
  the relation \((\exists x)((x, y) \in \mathbf{G})\) is equivalent to \(y \in \mathbf{B}\) .
\end{enumerate}

For it is sufficient to take A (resp. B) to be the set of all objects of the form \(pr_{1}z\) (resp. \(pr_{2}z\) ), where \(z \in G\) (§ 1, no. 6). Precisely,

\[
\mathrm{A} = \mathcal {E} _ {x} ((\exists y) ((x, y) \in \mathrm{G}))) \quad \text { and } \quad \mathrm{B} = \mathcal {E} _ {y} ((\exists x) ((x, y) \in \mathrm{G}));
\]

these sets are called respectively the first and second projections of the graph G, or the domain and range of G; they are denoted by \(pr_{1}\langle G\rangle\) and \(pr_{2}\langle G\rangle\) (or by \(pr_{1}G\) and \(pr_{2}G\) if there is no risk of confusion). It is immediately verified that \(G\subset(\mathrm{pr}_{1}\mathrm{G})\times(\mathrm{pr}_{2}\mathrm{G})\) ; every set of ordered pairs is therefore a subset of a product, and conversely. If one of the two sets \(pr_{1}G\) , \(pr_{2}G\) is empty, we have \(G=\varnothing\) (§2, Proposition 2).

Remark. The relation \(x = y\) has no graph, for the first projection of the graph, if it existed, would be the set of all objects (cf.~§1, no. 7, Remark).

\textbf{DEFINITION 2.} A correspondence between a set A and a set B is a triple

\[
\Gamma = (\mathrm{G}, \mathrm{A}, \mathrm{B}),
\]

where G is a graph such that \(pr_{1}G \subset A\) and \(pr_{2}G \subset B\) . G is said to be the graph of \(\Gamma\) , A is the source, and B the target of \(\Gamma\) .

If \((x, y) \in G\) , we say that ``y corresponds to x in the correspondence \(\Gamma\) ''. For each \(x \in pr_{1}G\) the correspondence \(\Gamma\) is said to be defined at x, and \(pr_{1}G\) is called the domain of \(\Gamma\) ; each \(y \in pr_{2}G\) is said to be a value taken by \(\Gamma\) , and \(pr_{2}G\) is called the range of \(\Gamma\) .

If \(R\{x, y\}\) is a relation which has a graph \(G\) (with respect to the letters \(x, y\) ), and if \(A, B\) are two sets such that \(\operatorname{pr}_1 G \subset A\) and \(\operatorname{pr}_2 G \subset B\) , we say that R is a relation between an element of A and an element of B (with respect to the letters x, y). The correspondence \(\Gamma = (G, A, B)\) is said to be the correspondence between A and B defined by the relation R (with respect to x and y).

Let \(\mathbf{G}\) be a graph and let \(\mathbf{X}\) be a set. The relation

\[
x \in \mathrm{X} \text {   and   } (x, y) \in \mathrm{G}
\]

implies \((x, y) \in G\) and therefore has a graph \(G'\) . The second projection of \(G'\) consists of all the objects which correspond under G to objects of X.

DEFINITION 3. Let G be a graph and X a set. The set of all objects which correspond under G to elements of X is called the image of X under G and is denoted by G⟨X⟩ or G(X).

Let \(\Gamma = (\mathbf{G},\mathbf{A},\mathbf{B})\) be a correspondence and let \(\mathbf{X}\) be a subset of \(\mathbf{A}\) . The set \(\mathbf{G}\langle \mathbf{X}\rangle\) is also denoted by \(\Gamma \langle \mathbf{X}\rangle\) or \(\Gamma (\mathbf{X})\) and is called the image of \(\mathbf{X}\) under \(\Gamma\) .

Remarks

\begin{enumerate}
\def\labelenumi{(\arabic{enumi})}
\item
  Precisely, \(\mathbf{G}\langle \mathbf{X}\rangle\) denotes the set \(\mathcal{E}_{\gamma}((\exists x)(x\in \mathbf{X}\) and \((x,y)\in \mathbf{G}))\) . From now on we shall not usually translate our definitions into formal language.
\item
  The notations G(X) and \(\Gamma(X)\) can occasionally lead to confusion with the notation introduced later (cf.~no. 4, Remark following Definition 9).
\end{enumerate}

Let G be a graph. Since the relation \((x, y) \in G\) implies \(y \in pr_{2}G\) , we have \(G\langle X \rangle \subset pr_{2}G\) for every set X. Since \((x, y) \in G\) implies \(x \in pr_{1}G\) , we have \(G\langle pr_{1}G \rangle = pr_{2}G\) . We have \(G\langle \phi \rangle = \phi\) , since \(x \notin \phi\) is a theorem. If \(X \subset pr_{1}G\) and \(X \neq \phi\) , we have \(G\langle X \rangle \neq \phi\) .

PROPOSITION 2. Let G be a graph and let X, Y be two sets; then the relation \(X \subset Y\) implies \(G\langle X\rangle \subset G\langle Y\rangle\) .

This is an immediate consequence of the definitions and of C50 (§ 1, no. 4).

COROLLARY. If \(\mathbf{A} \supset \mathrm{pr}_1\mathbf{G}\) , we have \(\mathbf{G}\langle \mathbf{A} \rangle = \mathrm{pr}_2\mathbf{G}\) .

DEFINITION 4. Let G be a graph and x an object. The set G\textless\{x\}\textgreater(which is sometimes denoted by G(x), by abuse of language) is called the section of G at x.

It follows immediately from C43 (Chapter I, § 5, no. 1) that the relation \(y \in G \langle \{x\} \rangle\) is equivalent to \((x, y) \in G\) . If \(G\) and \(G'\) are two graphs, the relation \(G \subset G'\) is thus equivalent to

\[
(\forall x) (\mathrm{G} \langle \{x \} \rangle \subset \mathrm{G} ^ {\prime} \langle \{x \} \rangle).
\]

If \(\Gamma = (G, A, B)\) is a correspondence between \(A\) and \(B\) , then for every \(x \in A\) the section of \(G\) at \(x\) is also called the section of \(\Gamma\) at \(x\) and is denoted by \(\Gamma \langle \{x\} \rangle\) (or \(\Gamma(x)\) ).

